\section{Plan, Challenges and Next Steps}
\label{sec:plan}
% Progress Report only. Required: the plan (Gantt chart), what is done, the challenges and the next steps.

\subsection{Status}
% Update this table on the day of submission.
\begin{table}[h]
\centering
\begin{tabular}{@{}lp{9.5cm}@{}}
\toprule
Component & Status \\
\midrule
Landmark extraction (MediaPipe) & implemented in the recording tool \\
Recording, viewing and validation tools & done \\
Normalisation & done and tested \\
Temporal smoothing (1-euro filter) & done; parameters provisional (one hand) \\
Static data loader & done \\
Static classifier & first pilot on one hand; not a result yet \\
Data collection & pilot clips only; several users still to record \\
Gesture segmentation & not implemented \\
Dynamic classifier & not implemented \\
Calibration, false-trigger prevention & not implemented \\
Command mapping & postponed; target device not decided \\
\bottomrule
\end{tabular}
\caption{Status of the project. \todo{update before submission}}
\end{table}

\subsection{Schedule}
% DRAFT bars: adjust to the real plan and add the names of the people (weeks start on Mondays;
% final report 4 January, presentations 8 and 15 January).
\begin{figure}[h]
\centering
\resizebox{\textwidth}{!}{%
\begin{ganttchart}[hgrid, vgrid, x unit=0.9cm, y unit title=0.5cm, y unit chart=0.55cm,
  title/.append style={fill=gray!15}, bar/.append style={fill=blue!40},
  milestone/.append style={fill=red!60}]{1}{8}
  \gantttitle{Week starting on (2026--2027)}{8} \\
  \gantttitlelist{23 Nov,30 Nov,7 Dec,14 Dec,21 Dec,28 Dec,4 Jan,11 Jan}{1} \\
  \ganttbar{Data collection, several users}{1}{3} \\
  \ganttbar{Orientation-robust features, static classifier}{1}{3} \\
  \ganttbar{Segmentation}{2}{3} \\
  \ganttbar{Dynamic classifier}{3}{5} \\
  \ganttbar{Calibration, false-trigger prevention}{4}{5} \\
  \ganttbar{Live system}{5}{6} \\
  \ganttbar{Experiments E1--E9}{5}{6} \\
  \ganttbar{Report writing}{6}{7} \\
  \ganttmilestone{Final report (4 Jan)}{7} \\
  \ganttbar{Presentations (8 and 15 Jan)}{7}{8}
\end{ganttchart}}
\caption{Project schedule. \todo{draft: adjust the bars and add who is responsible}}
\end{figure}

\subsection{Challenges}
% FACTS found so far (turn the ones that matter into text):
%  - the first static classifier learned the TILT of the hand: only the thumb-up clips were tilted
%    in the training data, and when the same hand was recorded with a tilt in another session the
%    open palm was called thumb up in 22 % of the frames; the fix is varied orientations in the
%    recordings, rotation-independent features or rotation augmentation;
%  - open palm is both a static gesture and the hand shape of the swipes, so they are separated by
%    movement (speed gate, hold time, cooldown) and this has to be validated;
%  - after a swipe the hand comes back by itself, so a swipe is defined as an out-and-back cycle;
%  - recording runs at about 20 frames per second, which leaves only 10 to 20 frames for a swipe;
%  - the data are still from one hand; generalisation to other users is not known.
\todo{challenges}

\subsection{Next Steps}
% WRITE: record several users and the missing gestures, handle orientation, segmentation, dynamic
% classifier, calibration and false-trigger prevention, then experiments E1 to E9.
\todo{next steps}
